\documentclass{article}
\usepackage[utf8]{inputenc}
\usepackage{amsmath, amssymb, amsthm}
\usepackage{geometry}
\usepackage{hyperref}
\geometry{a4paper, margin=2.5cm}
\title{HW5}
\date{\today}
\begin{document}
\maketitle
## Porosity and Permeability Analysis Report

### Introduction
This report presents an analysis of porosity and permeability data obtained from core samples under atmospheric conditions and at a minimum confining stress of 800 psig.

### Data Overview
The dataset includes:
- Porosity measurements at atmospheric conditions
- Grain density measurements at atmospheric conditions
- Permeability measurements at 800 psig
- Porosity measurements at 800 psig

### Porosity vs Grain Density at Atmospheric Conditions
A scatter plot of porosity vs grain density at atmospheric conditions was generated. The plot shows a moderate negative correlation, indicating that as grain density increases, porosity tends to decrease.

### Permeability vs Porosity at 800 psi
A scatter plot of permeability vs porosity at 800 psi was generated. The plot shows a positive correlation, suggesting that higher porosity values are associated with higher permeability values.

### Porosity Comparison
A scatter plot of porosity at 800 psi vs porosity at atmospheric conditions was generated. The plot shows a strong positive correlation, indicating that porosity values at both conditions are similar.

### Conclusion
The analysis provides insights into the relationships between porosity, grain density, and permeability under different conditions. These relationships are crucial for understanding reservoir properties and optimizing field development strategies.
\end{document}
